\section*{Bab \ref{ch:b1:complex} --- Bilangan Kompleks}

\begin{solution}{exo:b1:complex:1}
$1 + \iu = \sqrt 2\, \eu^{\iu\pi/4}$;
$\;\sqrt 3 - \iu = 2\, \eu^{-\iu\pi/6}$;
$\;-5 = 5\, \eu^{\iu\pi}$;
\[
\frac{1 + \iu}{\sqrt 3 - \iu}
= \frac{\sqrt 2}{2}\, \eu^{\iu(\pi/4 + \pi/6)}
= \frac{\sqrt 2}{2}\, \eu^{5\iu\pi/12}.
\]
Menghitung hasil bagi yang sama secara aljabar (kalikan dengan \omterm{def:b1:complex:field}{konjugatnya}):
\[
\frac{(1 + \iu)(\sqrt 3 + \iu)}{4}
= \frac{(\sqrt 3 - 1) + \iu(\sqrt 3 + 1)}{4}.
\]
Dengan menyamakannya terhadap $\frac{\sqrt 2}{2}(\cos\frac{5\pi}{12} +
\iu\sin\frac{5\pi}{12})$:
\[
\cos\frac{5\pi}{12} = \frac{\sqrt 6 - \sqrt 2}{4},
\qquad
\sin\frac{5\pi}{12} = \frac{\sqrt 6 + \sqrt 2}{4}.
\]
\end{solution}

\begin{solution}{exo:b1:complex:2}
$(1+\iu)^2 = 2\iu$, jadi $(1+\iu)^{20} = (2\iu)^{10} = 2^{10}\,\iu^{10}
= 1024 \times (-1) = -1024$.

Dalam bentuk eksponensial $(1+\iu)^n = 2^{n/2}\, \eu^{\iu n\pi/4}$, yang real jika
dan hanya jika $\sin\frac{n\pi}{4} = 0$, yakni $4 \mid n$. Jadi $(1+\iu)^n
\in \R$ tepat untuk kelipatan $4$ (dengan nilai $(-4)^{n/4}$).
\end{solution}

\begin{solution}{exo:b1:complex:3}
$\abs{z - 2} = \abs{z + \iu}$: berjarak sama dari titik $2$ dan
$-\iu$ --- yaitu sumbu tegak lurus ruas garis yang menghubungkan $(2, 0)$
dan $(0, -1)$.

$\abs{z - 1} = 2$: lingkaran berpusat $1$ dan berjari-jari $2$.

$\frac{z-1}{z+1} \in \iu\R$: menurut
\cref{prop:b1:complex:geometry}~(3), titik $M(z)$, $A(1)$,
$B(-1)$ memenuhi: garis $MA$ dan $MB$ tegak lurus (atau $z = 1$, yang
di sana hasil baginya $0 \in \iu\R$). \omterm{def:b1:logic:sets}{Himpunannya} adalah lingkaran berdiameter
$[-1, 1]$ (yaitu lingkaran satuan), dikurangi titik $-1$ yang di sana hasil baginya
tak terdefinisi. \emph{Periksa lewat perhitungan:} $z = \eu^{\iu\theta}$ memberikan
$\frac{z-1}{z+1} = \frac{\eu^{\iu\theta/2}(\eu^{\iu\theta/2} -
\eu^{-\iu\theta/2})}{\eu^{\iu\theta/2}(\eu^{\iu\theta/2} +
\eu^{-\iu\theta/2})} = \iu\tan\frac\theta2 \in \iu\R$.
\end{solution}

\begin{solution}{exo:b1:complex:4}
Pelinearannya:
\[
\sin^4\theta
= \Bigl(\frac{\eu^{\iu\theta} - \eu^{-\iu\theta}}{2\iu}\Bigr)^{\!4}
= \frac{\eu^{4\iu\theta} - 4\eu^{2\iu\theta} + 6 - 4\eu^{-2\iu\theta}
+ \eu^{-4\iu\theta}}{16}
= \frac{\cos 4\theta - 4\cos 2\theta + 3}{8}.
\]
Penjabarannya: menurut de Moivre, $\cos 4\theta = \Re\bigl((c +
\iu s)^4\bigr) = c^4 - 6c^2 s^2 + s^4$ dengan $c = \cos\theta$, $s =
\sin\theta$; lalu substitusikan $s^2 = 1 - c^2$:
\[
\cos 4\theta = c^4 - 6c^2(1 - c^2) + (1 - c^2)^2
= 8c^4 - 8c^2 + 1 .
\]
\end{solution}

\begin{solution}{exo:b1:complex:5}
$8\iu = 8\,\eu^{\iu\pi/2}$, jadi akar pangkat tiganya adalah $2\,\eu^{\iu(\pi/6 +
2k\pi/3)}$, $k = 0, 1, 2$:
\[
z_0 = 2\eu^{\iu\pi/6} = \sqrt 3 + \iu,\quad
z_1 = 2\eu^{5\iu\pi/6} = -\sqrt 3 + \iu,\quad
z_2 = 2\eu^{3\iu\pi/2} = -2\iu :
\]
yang membentuk segitiga sama sisi pada lingkaran berjari-jari $2$.

$-4 = 4\eu^{\iu\pi}$, jadi $z^4 = -4$ mempunyai penyelesaian $\sqrt 2\,
\eu^{\iu(\pi/4 + k\pi/2)}$: $\;1 + \iu$, $-1 + \iu$, $-1 - \iu$,
$1 - \iu$. Dengan memasangkan akar yang \omterm{def:b1:complex:field}{konjugat}:
\[
X^4 + 4 = \bigl(X^2 - 2X + 2\bigr)\bigl(X^2 + 2X + 2\bigr),
\]
karena $(X - (1+\iu))(X - (1-\iu)) = X^2 - 2X + 2$ dan serupa untuk
pasangan yang lain. (Jabarkan untuk memeriksanya.)
\end{solution}

\begin{solution}{exo:b1:complex:6}
$C_n + \iu S_n = \sum_{k=0}^{n} \eu^{\iu k\theta} =
\frac{\eu^{\iu(n+1)\theta} - 1}{\eu^{\iu\theta} - 1}$ (jumlah geometri,
dengan rasio $\eu^{\iu\theta} \neq 1$). Pemfaktoran setengah sudut
(\cref{met:b1:complex:trig}~(4)) pada pembilang dan penyebutnya:
\[
\frac{2\iu\sin\frac{(n+1)\theta}{2}\, \eu^{\iu(n+1)\theta/2}}
{2\iu\sin\frac{\theta}{2}\, \eu^{\iu\theta/2}}
= \frac{\sin\frac{(n+1)\theta}{2}}{\sin\frac{\theta}{2}}\,
\eu^{\iu n\theta/2}.
\]
Dengan mengambil bagian real dan bagian imajinernya:
\[
C_n = \frac{\sin\frac{(n+1)\theta}{2}}{\sin\frac{\theta}{2}}
\cos\frac{n\theta}{2},
\qquad
S_n = \frac{\sin\frac{(n+1)\theta}{2}}{\sin\frac{\theta}{2}}
\sin\frac{n\theta}{2}.
\]
\end{solution}

\begin{solution}{exo:b1:complex:7}
Diskriminannya: $\Delta = (3 + 4\iu)^2 - 4(-1 + 5\iu) = 9 + 24\iu - 16 +
4 - 20\iu = -3 + 4\iu$. Akar kuadrat $-3 + 4\iu$: selesaikan $x^2 - y^2
= -3$, $2xy = 4$, $x^2 + y^2 = 5$; lalu $x^2 = 1$, $y = 2/x$, dengan tanda
yang sama: $\delta = \pm(1 + 2\iu)$. Jadi
\[
z = \frac{(3 + 4\iu) \pm (1 + 2\iu)}{2}
\in \{\, 2 + 3\iu,\; 1 + \iu \,\}.
\]
\emph{Periksa:} jumlahnya $= 3 + 4\iu$ dan hasil kalinya $(2+3\iu)(1+\iu) = -1 +
5\iu$, sesuai tuntutan koefisiennya.
\end{solution}

\begin{solution}{exo:b1:complex:8}
\begin{enumerate}
  \item \cref{prop:b1:complex:sumroots} dengan $n = 5$.
  \item $u + v = \omega + \omega^2 + \omega^3 + \omega^4 = -1$ menurut (1).
        Untuk hasil kalinya, jabarkan lalu reduksi pangkatnya modulo $5$:
        \[
        uv = (\omega + \omega^4)(\omega^2 + \omega^3)
        = \omega^3 + \omega^4 + \omega^6 + \omega^7
        = \omega^3 + \omega^4 + \omega + \omega^2 = -1 .
        \]
        Jadi $u$ dan $v$ berjumlah $-1$ dan berhasil kali $-1$: keduanya
        adalah akar $X^2 + X - 1$.
  \item $u = \omega + \conj\omega = 2\cos\frac{2\pi}{5} > 0$ (karena
        sudutnya lancip), dan akar positif $X^2 + X - 1$ adalah
        $\frac{-1 + \sqrt 5}{2}$. Jadi $\cos\frac{2\pi}{5} =
        \frac{\sqrt 5 - 1}{4}$.
\end{enumerate}
\end{solution}

\begin{solution}{exo:b1:complex:9}
Jabarkan kedua kuadratnya seperti pada bukti
\cref{prop:b1:complex:rules}~(4):
\[
\abs{z + w}^2 = \abs z^2 + \abs w^2 + 2\Re(z\conj w),
\qquad
\abs{z - w}^2 = \abs z^2 + \abs w^2 - 2\Re(z\conj w),
\]
lalu jumlahkan. Secara geometris, $\abs{z+w}$ dan $\abs{z-w}$ adalah panjang
kedua diagonal jajargenjang bertitik sudut $0, z, z+w, w$,
sedangkan $\abs z$ dan $\abs w$ adalah panjang sisinya: jumlah
kuadrat diagonalnya sama dengan jumlah kuadrat keempat
sisinya.
\end{solution}

\begin{solution}{exo:b1:complex:10}
Misalkan $j = \eu^{2\iu\pi/3}$, sehingga $j^2 + j + 1 = 0$ dan $\eu^{\iu\pi/3} =
-j^2$, $\eu^{-\iu\pi/3} = -j$. Segitiganya sama sisi langsung jika dan hanya jika
$c - a = -j^2 (b - a)$, dan tak langsung jika dan hanya jika $c - a = -j(b - a)$
(\cref{ex:b1:complex:equilateral}).

Tinjau $P = a + jb + j^2 c$ dan $Q = a + j^2 b + jc$. Dengan memakai $1 +
j^2 = -j$ dan $j^3 = 1$:
\[
c - a + j^2(b - a) = c + j^2 b - (1 + j^2)\,a = c + j^2 b + ja
= j\,(a + jb + j^2 c) = jP,
\]
jadi kasus langsungnya berbunyi $jP = 0 \iff P = 0$; perhitungan yang sama
dengan $j$ menggantikan $j^2$ memberikan $c - a + j(b - a) = j^2 Q$, sehingga
kasus tak langsungnya berbunyi $Q = 0$. Jadi: sama sisi (dengan orientasi apa pun)
$\iff PQ = 0$. Dengan menjabarkannya, dan memakai $j + j^2 = -1$:
\[
PQ = a^2 + b^2 + c^2 + (j + j^2)(ab + bc + ca)
= a^2 + b^2 + c^2 - (ab + bc + ca).
\]
Jadi segitiganya sama sisi jika dan hanya jika $a^2 + b^2 + c^2 = ab +
bc + ca$.
\end{solution}

\begin{solution}{exo:b1:complex:11}
Karena $\omega^k$, $k = 0, \dots, n-1$, tepat merupakan $n$ akar
$X^n - 1$ (\cref{thm:b1:complex:roots}), dan polinomialnya
monik:
\[
X^n - 1 = \prod_{k=0}^{n-1} (X - \omega^k)
= (X - 1) \prod_{k=1}^{n-1} (X - \omega^k).
\]
Dengan membaginya dengan $X - 1$: $\;1 + X + \dots + X^{n-1} = \prod_{k=1}^{n-1}
(X - \omega^k)$. Menghitung nilainya di $X = 1$ memberikan $P = n$.

Selanjutnya $1 - \omega^k = -\eu^{\iu k\pi/n}\bigl(\eu^{\iu k\pi/n} -
\eu^{-\iu k\pi/n}\bigr) = -2\iu\,\eu^{\iu k\pi/n} \sin\frac{k\pi}{n}$,
jadi dengan mengambil \omterm{def:b1:complex:field}{modulus} pada $P = n$ (setiap $\sin\frac{k\pi}n > 0$ untuk $1 \leq k
\leq n-1$):
\[
n = \abs P = \prod_{k=1}^{n-1} 2\sin\frac{k\pi}{n}
= 2^{n-1} \prod_{k=1}^{n-1} \sin\frac{k\pi}{n},
\qquad\text{sehingga}\qquad
\prod_{k=1}^{n-1} \sin\frac{k\pi}{n} = \frac{n}{2^{n-1}} .
\]
\end{solution}

\begin{solution}{exo:b1:complex:12}
$z = 1$ bukan penyelesaian (karena $2^n \neq 0$), jadi persamaannya
setara dengan $\bigl(\frac{z+1}{z-1}\bigr)^n = 1$, yakni
$\frac{z+1}{z-1} = \omega^k$ dengan $\omega = \eu^{2\iu\pi/n}$ dan
$k \in \{0, \dots, n-1\}$. Nilai $k = 0$ tersingkir (karena $z + 1 =
z - 1$ mustahil). Untuk $1 \leq k \leq n - 1$, menyelesaikan
$z + 1 = \omega^k(z - 1)$ memberikan $z(1 - \omega^k) = -1 - \omega^k$,
sehingga dengan $\varphi = \frac{2k\pi}{n}$ dan pemfaktoran setengah sudut
$1 + \eu^{\iu\varphi} = 2\cos\frac\varphi2\,
\eu^{\iu\varphi/2}$, $\eu^{\iu\varphi} - 1 = 2\iu\sin\frac\varphi2
\,\eu^{\iu\varphi/2}$:
\[
z_k = \frac{1 + \omega^k}{\omega^k - 1}
= \frac{2\cos\frac{k\pi}{n}}{2\iu\,\sin\frac{k\pi}{n}}
= -\iu\,\frac{\cos\frac{k\pi}{n}}{\sin\frac{k\pi}{n}} ,
\]
yang murni imajiner sesuai klaimnya. \omterm{def:b1:logic:map}{Pemetaan} $t \mapsto \cos t/\sin t$ bersifat
\omterm{def:b1:logic:inj}{injektif} pada $\intoo0\pi$ (karena ia turun tegas), sehingga
$n - 1$ nilai $z_k$ itu berbeda dua-dua: persamaan itu, yang berderajat
$n - 1$ setelah dijabarkan (suku $z^n$ saling hapus), mempunyai tepat
$n - 1$ penyelesaian ini.
\end{solution}

\begin{solution}{pb:b1:complex:1}
\textbf{1.} $j = \cos\frac{2\pi}3 + \iu\sin\frac{2\pi}3 = -\frac12
+ \iu\frac{\sqrt3}2$; $j^2 = \eu^{4\iu\pi/3} = -\frac12 -
\iu\frac{\sqrt3}2 = \conj j$; $j^3 = 1$; $1 + j + j^2 = 0$ (jumlah
\omterm{def:b1:complex:unity}{akar satuan} pangkat tiga, \cref{prop:b1:complex:sumroots}); $j^{-1}
= j^2$ (karena $j \cdot j^2 = 1$). Pada lingkaran satuan, $1$, $j$,
$j^2$ adalah titik sudut segitiga sama sisi langsung.

\textbf{2.} Rotasi berpusat $a$ dan bersudut $\theta$ menetapkan $a$
dan memutar setiap vektor yang keluar dari $a$ sebesar $\theta$: $r(z) - a =
\eu^{\iu\theta}(z - a)$, yakni $r(z) = a + \eu^{\iu\theta}(z - a)$.
Mengomposisikan $r(z) = a + \eu^{\iu\theta}(z-a)$ dan $r'(z) = a' +
\eu^{\iu\theta'}(z-a')$:
\[
r' \circ r\,(z) = \eu^{\iu(\theta + \theta')} z +
\text{konstanta},
\]
yaitu \omterm{def:b1:logic:map}{pemetaan} berbentuk $Az + B$ dengan $A = \eu^{\iu(\theta+\theta')}$ yang
bermodulus $1$. Menurut \cref{prop:b1:complex:similitude}, ia adalah rotasi
bersudut $\arg A = \theta + \theta'$ bila $A \neq 1$, dan
translasi bila $A = 1$, yakni bila $\theta + \theta' \in 2\pi\Z$.

\textbf{3.} $(b, c, p)$ sama sisi jika dan hanya jika $\abs{p - b} = \abs{c -
b} = \abs{p - c}$. Dengan menulis $q = \frac{p - b}{c - b}$, kesamaan yang
pertama mengatakan $\abs q = 1$ dan yang kedua $\abs{q - 1} = 1$;
bersama-sama, $q = \eu^{\pm\iu\pi/3}$ (kedua titik potong
lingkaran $\abs q = 1$ dan $\abs{q - 1} = 1$ adalah $\frac12 \pm
\iu\frac{\sqrt3}2$). Jadi $p = b + \eu^{\pm\iu\pi/3}(c - b)$. Untuk
$a = 0$, $b = 1$, $c = \iu$ (sebuah segitiga langsung): pilihan
$\eu^{-\iu\pi/3}$ memberikan
\[
p = 1 + \Bigl(\tfrac12 - \iu\tfrac{\sqrt3}2\Bigr)(\iu - 1)
= \tfrac{1 + \sqrt3}2\,(1 + \iu) \approx 1.37 + 1.37\,\iu ,
\]
yang terletak di seberang garis $BC$ ($x + y = 1$) dari $a
= 0$: yaitu arah luar. Pilihan $\eu^{+\iu\pi/3}$ memberikan $p \approx -0.37
- 0.37\,\iu$, pada sisi yang sama dengan $a$: yaitu arah dalam.

\textbf{4.} Kalikan $P = a + jb + j^2c$ dengan $j^2$: $j^2 P = j^2 a +
j^3 b + j^4 c = b + jc + j^2 a$ (memakai $j^3 = 1$, $j^4 = j$); lalu
kalikan dengan $j$: $jP = ja + j^2 b + c$. Itulah kedua kesamaannya. Jika $P
= 0$ maka $j^2 P = jP = 0$: kriterianya berlaku untuk $(b, c, a)$ dan
$(c, a, b)$ juga --- yaitu keinvarianan siklik (memang harus demikian: sebuah
segitiga sama sisi tak peduli titik sudut mana yang didaftar lebih dulu).

\textbf{5.} Jika $a = b$: $0 = a(1 + j) + j^2 c = -j^2 a + j^2 c$
(memakai $1 + j = -j^2$), sehingga $c = a$. Jika $b = c$: $0 = a + b(j + j^2)
= a - b$, sehingga $a = b$. Jika $a = c$: $0 = a(1 + j^2) + jb = -ja + jb$,
sehingga $a = b$. Pada masing-masing kasus ketiga titiknya berimpit.

\textbf{6.} $(a'z + b') \circ (az + b) = a'a\,z + (a'b + b')$ dengan
$a'a \neq 0$: bentuknya sama. Invers $z \mapsto az + b$ adalah $z
\mapsto \frac1a z - \frac ba$, yang lagi-lagi berbentuk sama. Dengan
\omterm{def:b1:logic:map}{pemetaan} identitas sebagai unsur netralnya, semua kesebangunan langsung
membentuk grup terhadap komposisi.

\textbf{7.} $f(z) = az + b$ memenuhi $f(z_1) = w_1$, $f(z_2) =
w_2$ jika dan hanya jika $a z_1 + b = w_1$ dan $a z_2 + b = w_2$; dengan mengurangkannya,
$a(z_2 - z_1) = w_2 - w_1$, jadi
\[
a = \frac{w_2 - w_1}{z_2 - z_1} \;(\neq 0), \qquad
b = w_1 - a z_1
\]
terpaksa berlaku, dan sebaliknya pilihan itu memang memenuhi: itulah keberadaan
dan ketunggalannya.

\textbf{8.} $\abs{f(z) - f(w)} = \abs{a(z - w)} = \abs a\,\abs{z -
w}$: semua jarak dikalikan $\abs a$. Untuk bentuknya:
\[
\frac{f(c') - f(a')}{f(b') - f(a')} = \frac{a(c' -
a')}{a(b' - a')} = \frac{c' - a'}{b' - a'} .
\]
Jika dua segitiga $(a', b', c')$ dan $(a'', b'', c'')$ mempunyai bentuk
yang sama, misalkan $f$ satu-satunya kesebangunan langsung dengan $f(a') = a''$
dan $f(b') = b''$ (pertanyaan 7); maka bentuk $(a'', b'',
f(c'))$ sama dengan bentuk $(a', b', c')$, jadi juga bentuk $(a'', b'',
c'')$, dan bentuk itu menentukan titik sudut ketiga dari kedua titik sudut
pertamanya: $f(c') = c''$. Sebaliknya kesamaan bentuk menyusul dari
keinvarianan yang ditampilkan tadi.

\textbf{9.} $b = f(0) = 1$; $a + b = f(1) = \iu$ memberikan $a = \iu -
1$. Rasionya $\abs a = \sqrt2$, sudutnya $\arg(\iu - 1) =
\frac{3\pi}4$. Titik tetapnya:
\[
\zeta = \frac{b}{1 - a} = \frac1{2 - \iu} = \frac{2 + \iu}5 .
\]
Jadi $f$ adalah kesebangunan langsung berpusat $\frac{2+\iu}5$, berasio
$\sqrt2$ dan bersudut $\frac{3\pi}4$.

\textbf{10.} Dengan $\alpha + \beta = 1$ dan $f(z) = az + b$:
\[
f(\alpha a' + \beta b') = a\alpha a' + a\beta b' + b
= \alpha(a a' + b) + \beta(a b' + b)
= \alpha f(a') + \beta f(b') ,
\]
dengan kuncinya adalah $b = (\alpha + \beta)b$. Titik tengah ($\alpha = \beta =
\frac12$), titik berat (dengan mengiterasikannya), dan pusat Napoleon $\alpha b
+ \beta c$ di bawah karena itu bersifat \emph{ekuivarian}: mentransformasikan
segitiganya berarti mentransformasikan mereka juga. Jadi membuktikan sebuah
\omterm{def:b1:logic:statement}{pernyataan} yang invarian terhadap kesebangunan untuk satu segitiga yang ditempatkan baik
sudah membuktikannya untuk semua --- itulah izin yang dipakai pada pertanyaan 21.

\textbf{11.} Puncak arah luar pada $[b, c]$ adalah $p_a = b +
\eu^{-\iu\pi/3}(c - b)$ (pertanyaan 3), jadi pusatnya adalah
\[
n_a = \frac{b + c + p_a}3
= \frac{2b + c + \frac{1 - \iu\sqrt3}2\,(c - b)}3
= \frac{(3 + \iu\sqrt3)\,b + (3 - \iu\sqrt3)\,c}6
= \alpha b + \beta c ,
\]
dengan $\alpha = \frac{3 + \iu\sqrt3}6$ dan $\beta = \frac{3 -
\iu\sqrt3}6 = \conj\alpha$. Perhitungan yang sama pada sisi
$[c, a]$ dan $[a, b]$ memberikan $n_b = \alpha c + \beta a$ dan $n_c =
\alpha a + \beta b$ (yaitu pergeseran siklik perannya).

\textbf{12.} $\alpha + \beta = \frac{3 + \iu\sqrt3 + 3 -
\iu\sqrt3}6 = 1$. Jadi
\[
\frac{n_a + n_b + n_c}3
= \frac{(\alpha + \beta)(a + b + c)}3 = \frac{a + b + c}3 :
\]
kedua segitiga itu berbagi titik beratnya.

\textbf{13.} Kelompokkan menurut $\alpha$ dan $\beta$ lalu terapkan kesamaan
pertanyaan 4 pada $P = a + jb + j^2c$:
\[
n_a + j n_b + j^2 n_c
= \alpha\,(b + jc + j^2 a) + \beta\,(c + ja + j^2 b)
= \alpha\,j^2 P + \beta\,j P
= (\alpha j^2 + \beta j)\,P .
\]

\textbf{14.} Dengan $j = \frac{-1 + \iu\sqrt3}2$:
\[
\alpha j = \frac{(3 + \iu\sqrt3)(-1 + \iu\sqrt3)}{12}
= \frac{-3 + 3\iu\sqrt3 - \iu\sqrt3 - 3}{12}
= \frac{-3 + \iu\sqrt3}6 = -\beta ,
\]
sehingga $\alpha j + \beta = 0$, jadi $\alpha j^2 + \beta j = j(\alpha j
+ \beta) = 0$, dan pertanyaan 13 memberikan $n_a + jn_b + j^2n_c = 0$ untuk
setiap segitiga $(a, b, c)$. Menurut kriterianya (pertanyaan 4--5),
pusat-pusatnya membentuk segitiga sama sisi langsung atau sebuah titik tunggal:
itulah teorema Napoleon.

\textbf{15.} Puncak arah dalamnya adalah $b + \eu^{+\iu\pi/3}(c - b)$, dan
perhitungan pertanyaan 11 dengan $\eu^{+\iu\pi/3} = \frac{1 +
\iu\sqrt3}2$ menukarkan $\alpha$ dan $\beta$: $m_a = \beta b + \alpha
c$, $m_b = \beta c + \alpha a$, $m_c = \beta a + \alpha b$. Dengan memakai
kesamaan analognya $b + j^2c + ja = j\,Q$ dan $c + j^2a + jb
= j^2 Q$ untuk $Q = a + j^2b + jc$:
\[
m_a + j^2 m_b + j m_c
= \beta(b + j^2 c + ja) + \alpha(c + j^2 a + jb)
= (\beta j + \alpha j^2)\, Q = j(\beta + \alpha j)\,Q = 0 ,
\]
karena $\alpha j = -\beta$ (pertanyaan 14). Jadi pusat yang di dalam
memenuhi kriteria sama sisi \emph{tak langsung}: sama sisi dengan
orientasi yang berlawanan (atau sebuah titik).

\textbf{16.} Menurut pertanyaan 5 yang diterapkan pada tripel $(n_a, n_b, n_c)$
(yang memenuhi kriteria langsung), dua pusat berimpit jika dan hanya jika
ketiganya berimpit; dan ketiganya berimpit jika dan hanya jika \emph{kedua} kriterianya berlaku,
yakni jika dan hanya jika berlaku pula $n_a + j^2 n_b + j n_c = 0$. Hitunglah seperti pada
pertanyaan 13, dengan kesamaan $b + j^2c + ja = jQ$, $c + j^2a +
jb = j^2Q$:
\[
n_a + j^2 n_b + j n_c
= \alpha\,jQ + \beta\,j^2 Q = j(\alpha + \beta j)\,Q .
\]
Secara langsung: $\beta j = \frac{(3 -
\iu\sqrt3)(-1 + \iu\sqrt3)}{12} = \frac{-3 + 3\iu\sqrt3 +
\iu\sqrt3 + 3}{12} = \frac{\iu\sqrt3}3$, sehingga $\alpha + \beta j =
\frac{3 + \iu\sqrt3 + 2\iu\sqrt3}6 = \frac{1 + \iu\sqrt3}2 \neq
0$. Jadi segitiga Napoleon luarnya merosot jika dan hanya jika $Q = a +
j^2b + jc = 0$, yakni jika dan hanya jika $(a, b, c)$ merupakan segitiga sama sisi
tak langsung --- dan dalam hal itu semua konstruksi ``arah luar'' justru menunjuk
ke dalam daerah terlingkupi segitiganya dan berbagi satu pusat.

\textbf{17.} Dengan $a = 0$, $b = 1$, $c = \iu$:
\[
n_a = \alpha + \beta\iu = \frac{(3 + \sqrt3)(1 + \iu)}6,
\qquad
n_b = \alpha\iu = \frac{-\sqrt3 + 3\iu}6,
\qquad
n_c = \beta = \frac{3 - \iu\sqrt3}6 .
\]
Kuadrat sisinya: $n_a - n_b = \frac{(3 + 2\sqrt3) +
\iu\sqrt3}6$ memberikan $\abs{n_a - n_b}^2 = \frac{(3 + 2\sqrt3)^2 +
3}{36} = \frac{24 + 12\sqrt3}{36} = \frac{2 + \sqrt3}3$; $n_b -
n_c = \frac{(3 + \sqrt3)(-1 + \iu)}6$ memberikan $\abs{n_b - n_c}^2 =
\frac{2(3 + \sqrt3)^2}{36} = \frac{24 + 12\sqrt3}{36}$; $n_c - n_a
= \frac{-\sqrt3 - \iu(3 + 2\sqrt3)}6$ memberikan nilai yang sama lagi.
Ketiga kuadrat sisinya sama dengan $\frac{2 + \sqrt3}3$: jadi sama sisi,
sesuai janjinya.

\textbf{18.} Jabarkan:
\[
(a - b)(c - d) + (a - d)(b - c)
= (ac - ad - bc + bd) + (ab - ac - bd + cd)
= ab - ad - bc + cd ,
\]
dan $(a - c)(b - d) = ab - ad - bc + cd$: keduanya sama.

\textbf{19.} Ambil \omterm{def:b1:complex:field}{modulus} pada pertanyaan 18 lalu terapkan ketaksamaan
segitiga (\cref{prop:b1:complex:rules}~(4)):
\[
AC \cdot BD = \abs{(a-b)(c-d) + (a-d)(b-c)}
\leq \abs{a-b}\,\abs{c-d} + \abs{a-d}\,\abs{b-c}
= AB \cdot CD + AD \cdot BC ,
\]
dengan kesamaan jika dan hanya jika kedua sukunya terletak pada satu sinar garis yang sama dari $0$
(yakni kasus kesamaan pada ketaksamaan segitiga).

\textbf{20.} Pemfaktoran setengah sudut
(\cref{met:b1:complex:trig}~(4)):
$\eu^{\iu\theta} - \eu^{\iu\varphi} = 2\iu\,\sin\frac{\theta -
\varphi}2\;\eu^{\iu(\theta + \varphi)/2}$, dan mengambil modulusnya mematikan
faktor yang bermodulus satu: $\abs{\eu^{\iu\theta} - \eu^{\iu\varphi}}
= 2\,\abs{\sin\frac{\theta - \varphi}2}$.

\textbf{21.} Dengan $p = \eu^{\iu\theta}$, $\theta \in
\intoo{2\pi/3}{4\pi/3}$, pertanyaan 20 memberikan
\[
PA = 2\,\abs{\sin\tfrac\theta2},
\quad
PB = 2\,\abs{\sin\bigl(\tfrac\theta2 - \tfrac\pi3\bigr)},
\quad
PC = 2\,\abs{\sin\bigl(\tfrac\theta2 - \tfrac{2\pi}3\bigr)} .
\]
Pada busur itu, $\frac\theta2 \in \intoo{\pi/3}{2\pi/3}$: maka
$\sin\frac\theta2 > 0$; lalu $\frac\theta2 - \frac\pi3 \in
\intoo0{\pi/3}$, jadi sinus kedua bernilai positif; dan $\frac\theta2 -
\frac{2\pi}3 \in \intoo{-\pi/3}0$, jadi yang ketiga negatif sehingga
$PC = 2\sin\bigl(\frac{2\pi}3 - \frac\theta2\bigr)$. Dari jumlah ke
hasil kali:
\[
\sin\Bigl(\frac\theta2 - \frac\pi3\Bigr) +
\sin\Bigl(\frac{2\pi}3 - \frac\theta2\Bigr)
= 2\,\sin\frac\pi6\,\cos\Bigl(\frac\theta2 - \frac\pi2\Bigr)
= \sin\frac\theta2 ,
\]
jadi $PB + PC = 2\sin\frac\theta2 = PA$: itulah teorema van Schooten.

\textbf{22.} Pada $p = -1$: $PA = 2$, $PB = PC = 1$ dan semua sisi
segitiga sama sisinya berpanjang $\sqrt3$, sehingga kedua ruas
Ptolemeus terbaca $2\sqrt3$. Secara aljabar, dengan memakai $-1 - j^2 = j$ dan
$1 + j = -j^2$:
\[
(a - b)(p - c) = (1 - j)\,(-1 - j^2) = (1 - j)j = j - j^2
= \iu\sqrt3 ,
\]
\[
(a - c)(b - p) = (1 - j^2)(j + 1) = 1 + j - j^2 - j^3
= j - j^2 = \iu\sqrt3 .
\]
Kedua sukunya sama, jadi rasionya
$1 \in \R_{>0}$: yaitu kasus kesamaan pertanyaan 19, yang cocok persis dengan $PA
\cdot BC = PB \cdot AC + PC \cdot AB$.

\textbf{23.} Untuk $(a, b, c) = (1, j, j^2)$: $n_a = \alpha j +
\beta j^2 = -\beta + \beta j^2$ (pertanyaan 14) $= \beta(j^2 - 1)$;
secara numerik $\beta(j^2 - 1) = \frac{(3 - \iu\sqrt3)}6 \cdot
\bigl(-\frac32 - \iu\frac{\sqrt3}2\bigr) = -1$. Demikian pula $n_b =
\alpha j^2 + \beta = -j$ dan $n_c = \alpha + \beta j = -j^2$. Jadi
segitiga Napoleon luar dari $(1, j, j^2)$ adalah $(-1, -j, -j^2)$: yaitu
segitiga aslinya yang dicerminkan lewat titik beratnya $0$ --- berukuran
sama, terputar setengah putaran. Segitiga sama sisi adalah bentuk tetap
bagi konstruksi Napoleon, bukan limit yang menyusut.

\textbf{24.} (i) Perkalian dengan bilangan bermodulus satu sebagai rotasi
membangun puncak dan pusatnya (pertanyaan 2--3, 11) dan mengubah
jarak pada lingkaran menjadi sinus (pertanyaan 20). (ii) Struktur
grupnya beserta bentuk yang invarian membenarkan penormalan
lingkaran luarnya menjadi lingkaran satuan dan segitiganya menjadi $(1, j, j^2)$
pada pertanyaan 21, lewat keekuivarianan pertanyaan 10. (iii) Pemfaktoran
setengah sudut menggerakkan baik pertanyaan 20 maupun langkah
jumlah-ke-hasil-kali yang menuntaskan van Schooten.

\textbf{25.} Kamus itu mengubah \omterm{def:b1:logic:statement}{pernyataan} geometris menjadi
kesamaan polinomial yang setiap hipotesisnya berupa persamaan:
yang dibelinya adalah mekanisasi --- teorema Napoleon menyusut menjadi
$\alpha j + \beta = 0$, satu baris aritmetika di
$\Q(\iu\sqrt3)$; yang dibayarnya adalah keterlihatan geometris ---
perhitungan itu mengesahkan tetapi tidak memperlihatkan \emph{mengapa} pusat-pusatnya
menutup menjadi segitiga sama sisi. Jawaban yang adil adalah bahwa
kesamaan itu menjelaskan \emph{keniscayaan} teoremanya (ia berlaku
secara identik dalam $a, b, c$, jadi tak ada kecerdikan konfigurasi yang
terlibat), sedangkan gambarnya menjelaskan \emph{isinya}. Kamus yang sama
kembali dalam bentuk matriks ketika rotasi ini menjelma menjadi
matriks ortogonal $2 \times 2$ pada \cref{ch:b1:matrices} dan
\cref{ch:b1:euclid}, yang di sana ``perkalian dengan
$\eu^{\iu\theta}$'' menjadi prototipe isometri linear.

\end{solution}
